\section{Sharp support in projective extraspecial normalizers}
\label{sec:extraspecial-sharp-support}

\paragraph{Scope and provenance.}
D-TD08--10 refine D-TD05--07 separately.  The parent proposed investigating
exact support and the small-prime boundary; the proofs below are the
author's subsequent derivation. The author, root, and independent reviewer have each completed a whole-proof reading with PASS and no mathematical repair; the independent exact matrix supplement also passed.  They concern actual projective tensor factors and arbitrary
tensor partners.  They do not assume twelve nonidentity projected marks
and do not give a bound on the extraspecial exponent \(m\).

\paragraph{Notation and the proved interface.}
Let \(Q=\mathbf F_p^{\,2m}\), \(m\ge1\), act projectively irreducibly on
a complex vector space \(W\), with nondegenerate commutator pairing
\[
 L_vL_wL_v^{-1}L_w^{-1}
       =\zeta^{\langle v,w\rangle}I,\qquad a=\dim W=p^m.
 \tag{TDS.1}
\]
Here \(\zeta\) is a primitive \(p\)-th root and the pairing is a
nondegenerate alternating form over \(\mathbf F_p\).  For a normalizer
\(A\), let \(g\in\operatorname{Sp}(Q)\) be its action on \(Q\), and put
\[
 r=\operatorname{rank}(g-I),\qquad
 \nu(A)=a-\max_\lambda\dim\ker(A-\lambda I).
 \tag{TDS.2}
\]
Scalars do not affect \(\nu\), and representatives may be unitarized.
A finite projective subgroup has a finite inverse image in
\(\operatorname{SL}(W)\); alternatively a normalizer of the unitary
Weyl lifts has \(A^\dagger A\) scalar by irreducibility.
We use exactly the conclusions of D-TD05:
\[
 \begin{split}
 &\{L_v:v\in Q\}\text{ is a basis of }\operatorname{End}(W),\qquad
   \operatorname{tr}L_v=0\quad(v\ne0),\\
 &C_{\operatorname{PGL}(W)}(Q)=Q,\qquad
 A L_v A^{-1}=\eta_vL_{gv},\\
 &\eta|_{\operatorname{Fix}(g)}\text{ is a character},\qquad
 |\operatorname{tr}A|^2
    =\sum_{v\in\operatorname{Fix}(g)}\eta_v
    \in\{0,p^{\,2m-r}\}.
 \end{split}
 \tag{TDS.3}
\]
The character assertion is only on the fixed space; no global linear
phase convention is assumed.  For \(v\ne0\), scale \(L_v\) to order
\(p\).  Its proper powers have trace zero, so Fourier inversion gives
each of its \(p\) eigenvalues multiplicity \(a/p\), as in D-TD06.

\subsection{D-TD08: the sharp universal floor and tensor partners}

\begin{theorem}[D-TD08]
Every nonscalar normalizer satisfies
\[
                    \frac{\nu(A)}a\ge\frac{p-1}{2p}\ge\frac14.
 \tag{TDS.4}
\]
For every odd prime \(p\) and every \(m\ge1\), the first bound is
attained in the full projective normalizer.  At \(p=2\) it is attained
for every \(m\ge2\); for \(m=1\) the exact minimum is \(1/2\).
These are optimality and availability statements, not a classification
of all arbitrary-order equality elements or a claim that a specified
source tuple contains them.

The same normalized bound applies to the moving rank of \(A\otimes B\)
for every finite-order tensor partner \(B\), with arbitrary scalars.
In addition, when \(r>0\),
\[
 \frac{\nu(A)}a\ge\frac{1-p^{-r/2}}2;
 \qquad \operatorname{tr}A=0\ \Longrightarrow\ \frac{\nu(A)}a\ge\frac12.
 \tag{TDS.5}
\]
\end{theorem}

\begin{proof}
If \(g=I\), (TDS.3) makes \(A\) scalar times \(L_u\), \(u\ne0\),
whose equal eigenvalue multiplicities give \(\nu(A)/a=1-1/p\).
If \(g\ne I\), choose \(v\) with \(gv\ne v\).  Let \(E\) be a largest
eigenspace of \(A\), with eigenvalue \(\lambda\).  Every
\(z\in E\cap L_vE\) is fixed by \(C=A L_v A^{-1}L_v^{-1}\):
since \(L_v^{-1}z\in E\), direct substitution gives \(Cz=z\).
The operator \(C\) is scalar times the nonidentity Weyl operator
\(L_{gv-v}\), so its \(1\)-eigenspace has dimension at most \(a/p\).
Consequently
\[
                   a-2\nu(A)\le\dim(E\cap L_vE)\le a/p.
 \tag{TDS.6}
\]
This proves (TDS.4), without a trace estimate or normalizer splitting.
For (TDS.5), scale the largest eigenvalue to one.  The other eigenvalues
have absolute value one, so \(|\operatorname{tr}A|\ge a-2\nu(A)\);
apply (TDS.3).

For sharpness at odd \(p\), use the standard Weyl model on the basis
\(e_x\), \(x\in\mathbf F_p^m\).  Parity on the first coordinate,
\(e_{(x_1,x')}\mapsto e_{(-x_1,x')}\), normalizes the Weyl group
and has multiplicities \((p+1)p^{m-1}/2\) and \((p-1)p^{m-1}/2\).
At \(p=2,m\ge2\), the operator
\[
                         e_x\longmapsto(-1)^{x_1x_2}e_x
 \tag{TDS.7}
\]
normalizes the Weyl group and has multiplicities \(3a/4,a/4\).
At \(p=2,m=1\), nonscalarity in dimension two gives \(\nu\ge1\);
a nonidentity Weyl operator attains it.

Finally diagonalize \(B\).  Each of its eigenspaces contributes at
most \(a-\nu(A)\) times its dimension to any specified eigenspace of
\(A\otimes B\).  Summing proves the tensor assertion for every target
scalar, independently of the partner's field, degree, or image.
\end{proof}

\subsection{D-TD09: every rank-one symplectic-action coset}

\begin{theorem}[D-TD09]
Suppose \(\operatorname{rank}(g-I)=1\).  For odd \(p\), the exact
possible normalized supports in its full \(Q\)-coset are
\[
                              1-\frac2p,\qquad 1-\frac1p.
 \tag{TDS.8}
\]
The first occurs for exactly \(p\) of the \(p^{2m}\) projective lifts
in that coset, and the second for the others.  Both occur for every
\(m\ge1\).  At \(p=2\), the exact values and counts are
\[
 \begin{array}{c|c|c}
 \nu(A)/a&\text{number in the coset}&\text{availability}\\ \hline
 1/2&2^{2m-1}+2&m\ge1\\
 3/4&2^{2m-1}-2&m\ge2 .
 \end{array}
 \tag{TDS.9}
\]
In particular no rank-one action has support \(1/4\).
The counts hold inside every actual projective image containing \(Q\)
and one element inducing \(g\), since it contains the entire coset.
\end{theorem}

\begin{proof}
A rank-one symplectic map necessarily has the form
\[
              g(v)=v+c\langle w,v\rangle w\quad(w\ne0,\ c\ne0).
 \tag{TDS.10}
\]
Indeed write \(g-I=w\varphi\).  Symplecticity gives
\(\varphi(x)\langle w,y\rangle-\varphi(y)\langle w,x\rangle=0\),
since \(\langle w,w\rangle=0\).  Choose \(y\) with
\(\langle w,y\rangle\ne0\).  Thus \(\varphi\) is a nonzero multiple
of \(\langle w,\cdot\rangle\), and \((g-I)^2=0\).

For odd \(p\), choose \(P_w\) of order \(p\), with \(\zeta^j\)-eigenspace
\(E_j\) of dimension \(a/p\).  For \(\alpha=c/2\ne0\), define
\[
                          S_w|_{E_j}=\zeta^{\alpha j^2}I.
 \tag{TDS.11}
\]
The operator \(L_v\) sends \(E_j\) to \(E_{j+t}\), where
\(t=\langle w,v\rangle\).  Comparing phases shows that
\(S_wL_vS_w^{-1}\) is scalar times \(L_{v+2\alpha t w}\).
It induces \(g\); the centralizer identity makes every projective
lift uniquely \([L_uS_w]\).

If \(u=t w\), choose \(L_u=P_w^t\).  Its phases are
\(\zeta^{\alpha j^2+t j}\).  Completing the square shows that the
largest fibre of this quadratic is exactly two: every fibre has
at most two points, and a nonzero square value has two.  The largest
multiplicity is \(2a/p\); these are the \(p\) choices on
\(\mathbf F_pw\).
If \(u\in w^\perp\setminus\mathbf F_pw\), normalize the independent
commuting operators \(P_u,P_w\) to order \(p\).
All nonidentity elements of their group have trace zero, so every
joint eigenspace has dimension \(a/p^2\).  On its \((j,k)\)-space,
\(P_uS_w\) has eigenvalue \(\zeta^{k+\alpha j^2}\).
For any target and each \(j\), exactly one \(k\) contributes;
the multiplicity is \(a/p\).
If \(\langle w,u\rangle\ne0\), \(P_uS_w\) cycles the \(p\) spaces
\(E_j\).  Its \(p\)-th power is scalar, since the product of phases
uses all \(j\) exactly once and \(P_u^p=I\).
For each of the \(p\) roots of that scalar, the eigenvector equation
freely specifies a component in one \(E_j\) and determines the
others.  Every multiplicity is \(a/p\).
Such \(u\) exist even for \(m=1\), proving all odd-prime assertions.

For \(p=2\), scale each Weyl lift to a Hermitian involution \(P_v\).
Commuting labels have commuting lifts, and \(P_w\) has equal
\(+1,-1\) eigenspaces.  The operator
\[
                          S_w=(I+iP_w)/\sqrt2
 \tag{TDS.12}
\]
induces (TDS.10), with \(c=1\): it fixes commuting \(P_v\) and
sends anticommuting \(P_v\) to \(iP_wP_v\).
All projective lifts are \([P_uS_w]\).
For \(u=0,w\), its two distinct eigenvalues each have multiplicity
\(a/2\).  If \(\langle u,w\rangle=1\), then
\[
 (P_uS_w)^2=I,\qquad
 \operatorname{tr}(P_uS_w)
  =(\operatorname{tr}P_u+i\operatorname{tr}(P_uP_w))/\sqrt2=0,
 \tag{TDS.13}
\]
so its eigenvalues \(1,-1\) have equal multiplicities.
There are \(2^{2m-1}\) such \(u\).
For \(u\in w^\perp\setminus\langle w\rangle\), all four joint
eigenspaces of \(P_u,P_w\) have dimension \(a/4\).
The four eigenvalues \(\epsilon(1+i\delta)/\sqrt2\),
\(\epsilon,\delta=\pm1\), are distinct.
There are \(2^{2m-1}-2\) such \(u\), none for \(m=1\).
This proves (TDS.9) completely.
\end{proof}

\subsection{D-TD10: projective involutions and actual sources}

\begin{theorem}[D-TD10]
Let \([A]\) have exact projective order two, and scale \(A^2=I\).
For odd \(p\), \(g\ne I\), \(r\) is positive and even, and
\[
                         \frac{\nu(A)}a=\frac{1-p^{-r/2}}2.
 \tag{TDS.14}
\]
Every even \(r\), \(2\le r\le2m\), occurs in the full normalizer.
In every realized coset of a symplectic involution of rank \(r\),
there are exactly \(p^r\) projective involutions, all with this support.

For \(p=2\), put \(N=g-I\), \(R=\operatorname{im}N\), \(K=\ker N\).
When \(r>0\),
\[
 R=K^\perp\subset K,\quad r\le m,\qquad
 b_N(\bar v,\bar w)=\langle v,Nw\rangle\quad\text{on }Q/K
 \tag{TDS.15}
\]
is well-defined, symmetric, and nondegenerate.
There are exactly \(2^{2m-r}\) projective involutions in every
realized coset of \(g\).
If \(b_N\) is nonalternating, all have support \(1/2\).
If \(b_N\) is alternating, then \(r\) is even, and the counts are
\[
 \begin{array}{c|c}
 \nu(A)/a&\text{number in the coset}\\ \hline
 (1-2^{-r/2})/2&2^r\\
 1/2&2^{2m-r}-2^r .
 \end{array}
 \tag{TDS.16}
\]
The second case is absent when \(r=m\).
Nonalternating type occurs for every \(1\le r\le m\), and
alternating type for every positive even \(r\le m\), in the full
normalizer.  If \(r=0\), the coset is \(Q\): its identity is scalar,
and its other \(2^{2m}-1\) elements are projective involutions of
support \(1/2\).
\end{theorem}

\begin{proof}
For odd \(p\), the symplectic \(+1,-1\) eigenspaces of \(g\) are
orthogonal and nondegenerate, so \(r\) is even.  It is positive
since the projective kernel \(Q\) has odd exponent.
On \(K=\operatorname{Fix}(g)\), the character \(\eta\) has exponent
dividing \(p\), whereas \(A^2=I\) gives \(\eta^2=1\).  Thus \(\eta=1\).
By (TDS.3), \(|\operatorname{tr}A|=p^{m-r/2}\).
The two eigenvalue multiplicities give (TDS.14).

To count in a realized coset, start with any \(A_0\) inducing \(g\).
Its square is projectively \(L_h\).  Since \(A_0\) commutes with
its square, \(gh=h\).  The square of \(L_uA_0\) has projective
label \(h+(I+g)u\).
For odd \(p\), \(I+g\) is multiplication by two on the fixed space
and zero on the negative space.  Thus \((I+g)u=-h\) has exactly
\(p^r\) solutions, giving all projective involutions.
Parity on \(r/2\) Weyl coordinates realizes every allowable even \(r\).

For \(p=2\), symplecticity and \(g^2=I\) give
\(N^\dagger=N\) and \(N^2=0\).  Therefore \(R=K^\perp\subset K\).
The formula for \(b_N\) is independent of representatives and
symmetric; its radical before quotienting is exactly \(K\).
A nondegenerate alternating form has even dimension.

Existence and counts do not require a split normalizer.
Start with \(A_0\) inducing \(g\), and write \([A_0^2]=[L_h]\).
For \(v\in K\), conjugation sends the Hermitian involution \(P_v\)
to \(\eta_vP_v\), with \(\eta_v=\pm1\).
Thus \(A_0^2\) commutes with all \(P_v\), \(v\in K\), so
\(h\in K^\perp=R\).  The square of \(P_uA_0\) has projective
label \(h+Nu\).  The equation \(Nu=h\) has \(2^{2m-r}\) solutions,
giving precisely all projective involutions in the coset.

For such an involution write \(A P_vA^{-1}=\epsilon_vP_{gv}\),
\(\epsilon_v=\pm1\).  The equality \(A^2=I\) gives
\(\epsilon_v\epsilon_{gv}=1\).
Conjugating \(P_vP_{gv}\), which is scalar times \(P_{Nv}\), and
interchanging its two factors proves
\[
                       \eta_{Nv}=(-1)^{\langle v,Nv\rangle}.
 \tag{TDS.17}
\]
If \(b_N\) is nonalternating, this character is nontrivial on \(R\),
hence on \(K\).  Formula (TDS.3) gives zero trace and support \(1/2\).

If \(b_N\) is alternating, (TDS.17) says \(\eta|_R=1\).
Its restriction to \(K\) descends to \(K/R\).
The other involution lifts are \(P_uA\), \(u\in K\).
They change this character by
\(v\mapsto(-1)^{\langle u,v\rangle}\).
The alternating form on \(K/R\) is nondegenerate; hence every
character of \(K/R\) occurs, each for exactly \(|R|=2^r\) choices.
Exactly \(2^r\) lifts have trivial fixed-space character and
nonzero trace, of absolute value \(2^{m-r/2}\).
The other lifts have trace zero.  Their two eigenvalue
multiplicities prove (TDS.16).

For availability, multiply \(r\) independent commuting symplectic
transvections on \(r\) separate Weyl coordinates.  This gives rank
\(r\) and nonalternating \(b_N\) with diagonal entries one,
for every \(1\le r\le m\).
Each two-coordinate operator (TDS.7) gives rank two and alternating
\(b_N\) with matrix
\(\left(\begin{smallmatrix}0&1\\1&0\end{smallmatrix}\right)\).
Take \(r/2\) disjoint such pairs to realize every positive even
\(r\le m\).
The proved counts include the absence of zero trace for alternating
\(r=m\).  If \(r<m\), multiplying by a nonidentity Weyl operator
on an unused coordinate explicitly produces zero trace.
For \(r=0\), use (TDS.3) and the equal multiplicities of Weyl operators.
This covers every case.
\end{proof}

\paragraph{An all-power rank-profile bound.}
For exact projective order \(o\), scale \(A^o=I\).
For \(1\le j<o\), if \(g^j=I\), then \(A^j\) is a nonscalar Weyl
operator and has trace zero.  Other powers satisfy (TDS.3).
Fourier inversion yields
\[
 \frac{\nu(A)}a\ge
 1-\frac1o\left(1+
       \sum_{\substack{1\le j<o\\g^j\ne I}}
                 p^{-\operatorname{rank}(g^j-I)/2}\right).
 \tag{TDS.18}
\]
Powers with a nontrivial fixed-space phase character can also be
deleted from the sum.  This is a bound, not an exact arbitrary-order
spectral formula.

\paragraph{Actual tuples and the Scott interface.}
Return to \(V=W\otimes U\) with actual product-one generating inertia
\(T_i\), \(e=\dim V\), and total moving rank \(2e+6\).
Let \(x_i\) be its image in the actual quotient
\(P\subset\operatorname{PGL}(W)\), and let \(k\) count nonidentity images.
The arbitrary-partner conclusion gives
\[
                         k\,\frac{p-1}{2p}\le2+\frac6e.
 \tag{TDS.19}
\]
At \(p=2\), let \(k_0,k_1,k_{\ge2}\) count, among these nonidentity
images, symplectic ranks zero, one, and at least two.
The exact kernel and rank-one cases imply
\[
                         2(k_0+k_1)+k_{\ge2}\le8+\frac{24}e.
 \tag{TDS.20}
\]
For \(e>24\) the integer on the left is at most eight.
Identity projected images may still be genuine original inertia;
they have only been discarded from this lower bound.

There is also a legitimate rank-sum constraint.  Put
\(H=\langle g_i\rangle\subset\operatorname{Sp}(Q)\).
These images generate \(H\) and have product one.
Scott's lemma in its arbitrary-field formulation
\cite[Lemma2.3, author version p.~7]{TD:GuralnickMaroti2011} gives
\[
              \sum_i\operatorname{rank}(g_i-I)
                        \ge4m-2\dim Q^H.
 \tag{TDS.21}
\]
Its right side before simplification is
\(\dim Q+\dim[H,Q]-\dim Q^H\); self-duality gives
\([H,Q]^\perp=Q^H\).
The citation is to the explicit statement in that published research
paper; the original provenance is \cite{TD:Scott1977}.

If the actual projective factor \(P\) is primitive, then \(Q^H=0\).
A nonzero fixed \(v\) would make its order-\(p\) Weyl operator
projectively invariant under \(P\).  Its \(p\) equal-dimensional
eigenspaces would be permuted by \(P\), a nontrivial system of
imprimitivity.  Thus the rank sum is at least \(4m\) in the
primitive case.
This and (TDS.19)--(TDS.20) do not by themselves bound \(m\):
a few high-rank images can supply a large rank sum.
There is no inference of twelve projected marks, tuple realization,
balanced Hodge signature, or prescribed polarization.

\paragraph{Use in an entire partner.}
If a larger image permutes several such tensor-factor algebras,
a fixed factor on which the action is nonscalar has the floor \(1/4\).
On a moved cycle of length \(\ell\ge2\), tensor conjugation puts
the twist into one diagonalizable operator.
In its eigenbasis only the \(a\) constant basis words have orbit
length one; all other words have length at least two.
Each word orbit contributes at most one vector to any eigenspace.
Thus an eigenspace has dimension at most
\(a+(a^\ell-a)/2\le3a^\ell/4\), since \(a,\ell\ge2\).
Tensor convolution handles the other cycles and fixed factors.
This supplies the same floor for a nonscalar action on an entire
partner assembled from extraspecial factors of unbounded degree.
It supplies no floor for an additional unclassified partner factor;
the entire partner must be covered by the stated classes.
